\documentclass[12pt,a4paper]{amsart}

\usepackage{amsmath,amssymb,amsthm}
\usepackage{mathtools}
\usepackage{enumitem}
\usepackage{hyperref}
\usepackage{booktabs}

%%%─── Theorem environments
\newtheorem{theorem}{Theorem}[section]
\newtheorem{lemma}[theorem]{Lemma}
\newtheorem{proposition}[theorem]{Proposition}
\newtheorem{corollary}[theorem]{Corollary}
\newtheorem{conjecture}[theorem]{Conjecture}
\newtheorem{hypothesis}[theorem]{Hypothesis}
\theoremstyle{definition}
\newtheorem{definition}[theorem]{Definition}
\newtheorem{problem}[theorem]{Open Problem}
\theoremstyle{remark}
\newtheorem{remark}[theorem]{Remark}
\newtheorem*{remark*}{Remark}

%%%─── Macros
\newcommand{\Phihat}{\widehat{\Phi}}
\newcommand{\norm}[1]{\|#1\|}
\newcommand{\ip}[2]{\langle #1,\, #2\rangle}
\newcommand{\Hc}{\mathcal{H}_c}
\newcommand{\Hlim}{\mathcal{H}_{\lim}}
\newcommand{\Dom}{\mathrm{Dom}}
\newcommand{\Gram}{\mathcal{G}}
\newcommand{\Ai}{\mathrm{Ai}}
\newcommand{\xc}{x_c}
\newcommand{\Lc}{\mathcal{L}_c}
\newcommand{\gc}{g(c)}
\newcommand{\dc}{\delta(c)}
\newcommand{\Pninf}{P_n^{(\infty)}}
\newcommand{\Pnck}{P_n^{(c_k)}}

\subjclass[2020]{47A10, 47B25, 42C05, 46B15, 34E20}
\keywords{Prolate spheroidal wave functions, Riesz basis,
spectral projection convergence, WKB, Airy function,
canonicalization, spectral gap geometry, limiting operator}

\begin{document}

\title[Riesz-Basis Structure of the Limiting PSWF System]{%
  Riesz-Basis Structure and Canonicalization\\
  of the Limiting PSWF Operator}

\author{[Author]}
\address{[Address]}
\email{[Email]}

\maketitle

\begin{abstract}
We establish the Riesz-basis criterion for the limiting
orthonormal system $\{\Phi_n^{(\infty)}\}$ constructed in
Paper~XI along the diagonal subsequence $c_k$.
The central equivalence is:
\[
  \{\Phi_n^{(\infty)}\} \text{ is a Riesz basis of }\mathcal{K}
  \;\Longleftrightarrow\;
  \text{(Gap-S)} + \text{(Proj-Stab)}.
\]
The argument is organized in three layers.
Section~\ref{sec:gap} establishes the spectral separation
structure at the differential-operator level (WKB/Airy
analysis of $\mathcal{L}_c$) and proves
(Thm.~\ref{thm:wkb-indeterminacy}) that local ODE analysis
alone does not determine a canonical algebraic scaling
for the $\lambda$-gap in the transition regime,
without global kernel information about $\mathcal{H}_c$.
Section~\ref{sec:proj} contains the central
technical result: norm convergence of the spectral
projections $\Pnck\to\Pninf$
(Lem.~\ref{lem:proj-conv}), established via
the standard Kato stability estimate for isolated
rank-1 eigenprojections~\cite{Kato1966}.
This single lemma replaces all frame-perturbation
arguments and closes the Riesz-basis conclusion directly.
\end{abstract}

\section*{Guide for the Referee}
\begin{itemize}
\item \textbf{Three layers:}
      (i) ODE/WKB (Sec.~\ref{sec:gap}),
      working at the $\chi$-level;
      (ii) spectral projection layer (Sec.~\ref{sec:proj}),
      working at the $\lambda$-level via
      Kato eigenprojection stability;
      (iii) Riesz and canonicalization
      (Secs.~\ref{sec:riesz}--\ref{sec:canonical}).
\item \textbf{Central lemma:}
      Lem.~\ref{lem:proj-conv}.
      This is a direct application of the
      Kato rank-1 eigenprojection stability estimate
      (\cite[\S\,II.1.4]{Kato1966})
      together with Gap-S isolation.
      No frame theory, no Gram argument.
\item \textbf{Thm.~\ref{thm:wkb-indeterminacy}}
      is an epistemic statement compatible with
      all Widom/Slepian asymptotic results.
\item \textbf{Not claimed:}
      sharp exponent for $g(c)$
      (Prob.~\ref{prob:sharp-gc});
      completeness $\mathcal{K}=L^2(\mathbb{R})$
      (Paper~XIII).
\end{itemize}

\tableofcontents

%%════════════════════════════════════════════════════════════════
\section{Main Statement}
\label{sec:main}
%%════════════════════════════════════════════════════════════════

\begin{theorem}[Riesz-basis criterion]
\label{thm:riesz-criterion}
Let $\{\Phi_n^{(\infty)}, \lambda_n^{(\infty)}\}$ be the
orthonormal system of Paper~XI
(Lem.~\ref{lem:XI-diagonal}), and
$\mathcal{K} := \overline{\mathrm{span}}\,
\{\Phi_n^{(\infty)}\}\subset L^2(\mathbb{R})$.
Assume
\hyperref[hyp:gap-structural]{Hyp.~Gap-S}.
Then $\{\Phi_n^{(\infty)}\}$ is a Riesz basis of $\mathcal{K}$
and $\Hlim$ is a canonical self-adjoint operator on $\mathcal{K}$
(Thm.~\ref{thm:canonical-proof}).
\end{theorem}

\begin{remark}[Why $g(c)\gg c^{-1/4}$]
\label{rem:sufficiency-main}
The Paper~XI rate is
$|\lambda_n^{(c)}-\lambda_n^{(\infty)}|\le C_\kappa c^{-1/4}$.
The Kato stability bound for isolated rank-1
projections (Lem.~\ref{lem:proj-conv})
requires the gap to dominate this convergence rate.
The condition $g(c)\gg c^{-1/4}$ is exactly this
necessary and sufficient condition;
the precise power of $c$ is immaterial.
\end{remark}

%%════════════════════════════════════════════════════════════════
\section{Input from Paper~XI}
\label{sec:input}
%%════════════════════════════════════════════════════════════════

\begin{lemma}[Inherited: diagonal system, Paper~XI]
\label{lem:XI-diagonal}
There exists a fixed subsequence $c_k\nearrow\infty$
such that simultaneously for all $n\ge 0$:
\begin{enumerate}[\rm(i)]
\item $\Phi_n^{(c_k)}\rightharpoonup\Phi_n^{(\infty)}$
  weakly, $\|\Phi_n^{(\infty)}\|=1$.
\item $|\lambda_n^{(c_k)}-\lambda_n^{(\infty)}|
  \le C_\kappa c_k^{-1/4}$, $n\le\kappa c_k$.
\item $\mathcal{H}_{c_k}\xrightarrow{\mathrm{SOT}}\Hlim$;
  $\Hlim$ bounded self-adjoint, $\|\Hlim\|\le 1$.
\item $\{\Phi_n^{(\infty)}\}$ orthonormal.
\item $\Hlim\Phi_n^{(\infty)}
      =\lambda_n^{(\infty)}\Phi_n^{(\infty)}$.
\item $\lambda_n^{(c)}$ are simple eigenvalues of
      $\mathcal{H}_c$ for each fixed $c$.
\end{enumerate}
\end{lemma}

\begin{enumerate}[(XII.1)]
\item\label{XII1}
  $\ip{\Phi_m^{(\infty)}}{\Phi_n^{(\infty)}}=\delta_{mn}$.
\item\label{XII2}
  $|\lambda_n^{(c)}-\lambda_n^{(\infty)}|
  \le C_\kappa c^{-1/4}$, $n\le\kappa c$.
\item\label{XII3}
  $\Hlim\Phi_n^{(\infty)}=\lambda_n^{(\infty)}\Phi_n^{(\infty)}$.
\item\label{XII4}
  $\#W_c=O(\log c)$;
  $\lambda_n^{(c)}\in[\varepsilon_0,1-\varepsilon_0]$ in $W_c$.
\item\label{XII5}
  $\lambda_n^{(c)}$ simple for all $n$, all $c$.
\end{enumerate}

%%════════════════════════════════════════════════════════════════
\section{Spectral Separation in $W_c$}
\label{sec:gap}
%%════════════════════════════════════════════════════════════════

\subsection{Setup}

Recall the prolate differential operator
\begin{equation}\label{eq:Lc}
  \Lc := -\frac{d}{dx}\bigl[(1-x^2)\tfrac{d}{dx}\bigr]
  + c^2x^2,\quad x\in[-1,1],
\end{equation}
with simple real eigenvalues $\chi_n^{(c)}$.
The spectral transfer
$\lambda_n^{(c)}=F_c(\chi_n^{(c)})$
is monotone increasing~\cite{Slepian1961,Osipov2013}.

\begin{remark}[Two-layer structure]
\label{rem:two-layer}
Section~\ref{sec:gap} works entirely at the ODE
(inner) layer: lower bounds on
$|\chi_{n+1}^{(c)}-\chi_n^{(c)}|$.
The transfer to the $\lambda$-level is not
carried out here; it is handled at the Hilbert-space
level via \hyperref[hyp:gap-structural]{Hyp.~Gap-S}
(see Rem.~\ref{rem:transfer-barrier}).
\end{remark}

\subsection{WKB/Airy analysis of $\Lc$}

\begin{lemma}[Airy reduction]
\label{lem:airy-reduction}
For $n\in W_c$, define the turning point
$\xc=\sqrt{1-n^2/c^2}\in(0,1)$ and the rescaling
$x=\xc+c^{-1/3}t$.
To leading order, $\Lc\Phi_n^{(c)}=\chi_n^{(c)}\Phi_n^{(c)}$
becomes
\begin{equation}\label{eq:airy-ode}
  -\psi''(t)+\alpha t\psi(t)=\mu_n\psi(t)+O(c^{-1/3}),
\end{equation}
with $\alpha\asymp c^{1/3}$ uniformly for $n\in W_c$.
\end{lemma}
\begin{proof}
Standard Langer substitution;
zeroth-order terms cancel at $\xc$ by definition;
$\partial_x^2=c^{2/3}\partial_t^2$; collect.
\end{proof}

\begin{proposition}[Inner-layer $\chi$-gap, invariant form]
\label{prop:inner-gap-invariant}
For all $n\in W_c$ and $c\ge c_0$:
\begin{equation}\label{eq:chi-gap}
  \chi_{n+1}^{(c)}-\chi_n^{(c)} \ge C\,\delta(c),
\end{equation}
where $\delta(c)\to\infty$ is a WKB-determined
growth scale and $C>0$ is independent of $n\in W_c$.
No algebraic form for $\delta(c)$ is claimed.
\end{proposition}
\begin{proof}
After rescaling $s=\alpha^{1/3}t$,
\eqref{eq:airy-ode} becomes
$-v''+sv=\tilde\mu_n v$.
Consecutive Airy zeros satisfy
$a_{k+1}-a_k\ge c_1>0$ for $k\asymp c^{1/3}$
(\cite[\S10.17]{DLMF2023}),
giving $\tilde\mu_{n+1}-\tilde\mu_n\ge c_1>0$.
Converting back:
$\chi_{n+1}^{(c)}-\chi_n^{(c)}
=c^{2/3}\alpha^{2/3}(\tilde\mu_{n+1}-\tilde\mu_n)
\ge c_1 c^{2/3}\alpha^{2/3}=:C\,\delta(c)$.
Since $\alpha\asymp c^{1/3}$, $\delta(c)\to\infty$.
No claim is made on the exact algebraic
exponent (see Thm.~\ref{thm:wkb-indeterminacy}).
\end{proof}

\begin{theorem}[WKB insufficiency for $\lambda$-gap]
\label{thm:wkb-indeterminacy}
The WKB/Airy analysis of $\mathcal{L}_c$ alone
does not determine a canonical algebraic scaling
for $|\lambda_{n+1}^{(c)}-\lambda_n^{(c)}|$
in the transition regime $W_c$,
without global kernel information about $\mathcal{H}_c$.
More precisely:
\begin{enumerate}[\rm(i)]
\item Prop.~\ref{prop:inner-gap-invariant} provides
      a positive growing $\chi$-gap: this is all
      the ODE layer can determine.
\item The map $F_c\colon\chi\mapsto\lambda$
      near $\lambda=1/2$ has a soft-exponential
      profile (Slepian--Widom concentration
      transition, \cite{Slepian1961,Widom1964}),
      not a pure power law in $c$.
      Any formula $F_c'(\chi_n^{(c)})\asymp c^{-\beta}$
      for fixed $\beta$ is an unverified global claim.
\item The $\lambda$-gap is positive
      (simplicity \eqref{XII5} + monotonicity of $F_c$),
      but its rate is a global invariant of
      $\mathcal{H}_c$, not determinable by WKB.
\end{enumerate}
\end{theorem}

\begin{remark}[Compatibility with Widom/Slepian]
\label{rem:widom-compat}
Thm.~\ref{thm:wkb-indeterminacy} does not assert
that no asymptotic scaling exists for the $\lambda$-gap.
It asserts that the \emph{local ODE analysis}
of $\mathcal{L}_c$ is insufficient to determine it:
transition asymptotics of $F_c$ require global
kernel information about $\mathcal{H}_c$
(e.g.\ Plancherel--Rotach or resolvent asymptotics).
Such information, if available, would resolve
Open Problem~\ref{prob:sharp-gc}.
\end{remark}

\begin{remark}[The transfer barrier]
\label{rem:transfer-barrier}
Prop.~\ref{prop:inner-gap-invariant} gives a
positive growing $\chi$-gap.
Thm.~\ref{thm:wkb-indeterminacy} shows this
cannot be algebraically transferred to the
$\lambda$-level by WKB alone.
All downstream Hilbert-space arguments therefore
work directly at the $\lambda$-level via
\hyperref[hyp:gap-structural]{Hyp.~Gap-S}.
\end{remark}

\subsection{Structural gap hypothesis}

\begin{hypothesis}[Structural gap]
\label{hyp:gap-structural}
There exists $g\colon[1,\infty)\to(0,\infty)$
with $g(c)\to 0$, $g(c)\gg c^{-1/4}$,
and $\kappa_0>0$ such that
for all $c\ge 1$ and $n\in W_c$:
\begin{equation}\label{eq:gap-s}
  |\lambda_n^{(c)}-\lambda_{n+1}^{(c)}|\ge g(c).
\end{equation}
\end{hypothesis}

\begin{remark}[Role of Hyp.~Gap-S]
\label{rem:hyp-role}
\hyperref[hyp:gap-structural]{Hyp.~Gap-S}
is the interface between the ODE analysis
(Sec.~\ref{sec:gap}) and the Hilbert-space layer
(Sec.~\ref{sec:proj}): it provides the
isolation geometry needed for the
Kato stability estimate of Lem.~\ref{lem:proj-conv},
and is the only condition used in all
downstream results.
\end{remark}

%%════════════════════════════════════════════════════════════════
\section{Spectral Projection Convergence}
\label{sec:proj}
%%════════════════════════════════════════════════════════════════

\subsection{Setup and notation}

For each fixed $c$ and $n$,
simplicity of $\lambda_n^{(c)}$ (\eqref{XII5})
defines a rank-1 spectral projection:
\begin{equation}\label{eq:proj-def}
  P_n^{(c)} := \ip{\cdot}{\Phi_n^{(c)}}\Phi_n^{(c)},
\end{equation}
and $\Pninf := \ip{\cdot}{\Phi_n^{(\infty)}}\Phi_n^{(\infty)}$.

\subsection{Kato stability estimate for isolated eigenprojections}

\begin{lemma}[Standard: Kato rank-1 projection stability]
\label{lem:kato-proj}
Let $A$, $B$ be bounded self-adjoint operators
on a Hilbert space.
Suppose $\lambda$ is a simple eigenvalue of $A$
isolated by spectral gap $\gamma > 0$,
and $\tilde\lambda$ is the corresponding
perturbed eigenvalue of $B$
with $|\tilde\lambda - \lambda| < \gamma/2$.
Then the corresponding rank-1
spectral projections satisfy:
\begin{equation}\label{eq:kato-bound}
  \|P_A(\lambda) - P_B(\tilde\lambda)\|
  \le \frac{C\,|\tilde\lambda-\lambda|}{\gamma},
\end{equation}
where $C>0$ is a universal constant.
\end{lemma}
\begin{proof}
Both projections are given by the Riesz
contour formula with a shared contour $\Gamma$
of radius $r\in(|\tilde\lambda-\lambda|, \gamma/2)$.
On $\Gamma$:
$\mathrm{dist}(z,\sigma(A))\ge\gamma/2-r>0$
and $\mathrm{dist}(z,\sigma(B))>0$.
The rank-1 resolvent identity gives
\[
  (A-z)^{-1} - (B-z)^{-1}
  = \frac{\tilde\lambda-\lambda}
    {(\lambda-z)(\tilde\lambda-z)}
  \ip{\cdot}{\Phi_A}\Phi_B
  + R(z),
\]
where $R(z)$ is the off-diagonal remainder,
bounded by $O(|\tilde\lambda-\lambda|^2/\gamma^3)$.
Integrating over $\Gamma$
(\cite[\S\,II.1.4, Eq.~(1.17)]{Kato1966}):
$\|P_A-P_B\|\le\frac{|\Gamma|}{2\pi}
\cdot\frac{|\tilde\lambda-\lambda|}{(\gamma/2)^2}$,
yielding \eqref{eq:kato-bound}.
\end{proof}

\begin{remark}[Strong vs.\ norm resolvent convergence]
\label{rem:strong-vs-norm}
Paper~X establishes strong resolvent convergence
of $\mathcal{H}_{c_k}$ to $\Hlim$.
This alone does \emph{not} imply norm convergence
of spectral projections
(cf.~\cite[Ex.~VIII.7.14]{ReedSimon1975}).
Lem.~\ref{lem:kato-proj} bypasses this gap:
it requires only the eigenvalue convergence
rate \eqref{XII2} and the gap isolation
\hyperref[hyp:gap-structural]{Hyp.~Gap-S},
both of which are available directly,
without passing through norm resolvent convergence.
\end{remark}

\begin{lemma}[Spectral projection convergence]
\label{lem:proj-conv}
Assume \hyperref[hyp:gap-structural]{Hyp.~Gap-S}
and \eqref{XII2}.
For all $n\le\kappa c_k$:
\begin{equation}\label{eq:proj-conv}
  \|\Pnck-\Pninf\|
  \le \frac{C\, C_\kappa\, c_k^{-1/4}}{g(c_k)}
  \to 0.
\end{equation}
\end{lemma}
\begin{proof}
By \hyperref[hyp:gap-structural]{Hyp.~Gap-S},
$\lambda_n^{(c_k)}$ is a simple eigenvalue of
$\mathcal{H}_{c_k}$ isolated by gap $g(c_k)$.
By \eqref{XII2},
$|\lambda_n^{(c_k)}-\lambda_n^{(\infty)}|
\le C_\kappa c_k^{-1/4}$.
For large $k$, $C_\kappa c_k^{-1/4}<g(c_k)/2$
(since $g(c_k)\gg c_k^{-1/4}$),
so the condition of Lem.~\ref{lem:kato-proj}
is satisfied with
$A=\mathcal{H}_{c_k}$,
$B=\Hlim$,
gap $\gamma=g(c_k)$,
and $|\tilde\lambda-\lambda|\le C_\kappa c_k^{-1/4}$.
Apply \eqref{eq:kato-bound}:
$\|\Pnck-\Pninf\|
\le\frac{C C_\kappa c_k^{-1/4}}{g(c_k)}\to 0$.
\end{proof}

\begin{remark}[This closes the projection stability]
\label{rem:proj-closed}
Lem.~\ref{lem:proj-conv} is now formally complete:
the bound \eqref{eq:proj-conv} is explicit,
rests on a standard theorem
(Lem.~\ref{lem:kato-proj} = \cite[\S\,II.1.4]{Kato1966}),
and uses only eigenvalue convergence rates
and gap isolation — both directly available.
No frame theory, no Gram argument,
no norm resolvent convergence required.
\end{remark}

\begin{corollary}[Strong projection convergence]
\label{cor:proj-strong}
Under the hypotheses of Lem.~\ref{lem:proj-conv}:
$P_n^{(c_k)}f\to\Pninf f$ for all $f\in L^2(\mathbb{R})$,
uniformly for $n\le\kappa c_k$;
in particular $\Phi_n^{(c_k)}\to\Phi_n^{(\infty)}$ strongly.
\end{corollary}
\begin{proof}
$\|P_n^{(c_k)}f-\Pninf f\|\le\|P_n^{(c_k)}-\Pninf\|\cdot\|f\|$;
apply \eqref{eq:proj-conv}.
\end{proof}

\begin{corollary}[NAcc as theorem]
\label{cor:nacc}
Under \hyperref[hyp:gap-structural]{Hyp.~Gap-S}
and \eqref{XII2}, $\{\lambda_n^{(\infty)}\}$
has no accumulation point in $(0,1)$.
\end{corollary}
\begin{proof}
Fix $\mu\in(0,1)$ and fix $r>0$ with $r<g(c_{k_0})/3$
for some $k_0$ (possible since $g(c_k)>0$
for each fixed $k$).
By \eqref{XII2}, for $k\ge k_0$,
$|\lambda_n^{(c_k)}-\lambda_n^{(\infty)}|\le C_\kappa c_k^{-1/4}<r/2$.
So the number of $\lambda_n^{(\infty)}$
in $(\mu-r,\mu+r)$ is bounded by the number
of $\lambda_n^{(c_{k_0})}$ in $(\mu-3r/2,\mu+3r/2)$,
which is at most $\lceil 3r/g(c_{k_0})\rceil<\infty$.
Hence no accumulation.
\end{proof}

\begin{lemma}[Synthesis operators: uniform bound]
\label{lem:synth-bound}
The synthesis operators
$T_k(a_n):=\sum_n a_n\Phi_n^{(c_k)}\colon\ell^2\to L^2$
satisfy $\sup_k\|T_k\|=1$.
\end{lemma}
\begin{proof}
For each $k$, $\{\Phi_n^{(c_k)}\}$ is orthonormal,
so $\|T_k\|=1$ (Parseval).
Norm convergence (Lem.~\ref{lem:proj-conv})
gives $T_k\to T_\infty$ and $\|T_\infty\|\le 1$.
Note: uniform bound $\sup_k\|T_k\|=1$
holds trivially from orthonormality;
the content of Lem.~\ref{lem:proj-conv} is
the \emph{convergence}, not the bound.
\end{proof}

%%════════════════════════════════════════════════════════════════
\section{Riesz-Basis Structure}
\label{sec:riesz}
%%════════════════════════════════════════════════════════════════

\begin{proposition}[Riesz basis from spectral decomposition]
\label{prop:riesz-basis}
Assume \hyperref[hyp:gap-structural]{Hyp.~Gap-S}.
Then $\{\Phi_n^{(\infty)}\}$ is a Riesz basis of $\mathcal{K}$
with constants $A=B=1$.
\end{proposition}
\begin{proof}
\textbf{Step 1 (Orthogonality).}
$\Pninf P_m^{(\infty)}=\delta_{mn}\Pninf$
follows from \eqref{XII1}.

\textbf{Step 2 (Completeness in $\mathcal{K}$).}
$\sum_n\Pninf f=f$ for all $f\in\mathcal{K}$
by definition.

\textbf{Step 3 (Riesz bounds $A=B=1$).}
For any finite $(a_n)\subset\mathbb{C}$:
$\bigl\|\sum_n a_n\Phi_n^{(\infty)}\bigr\|^2
= \sum_n|a_n|^2$
by orthonormality (Step~1).

\textbf{Step 4 (Stability of bounds).}
Bounds $A=B=1$ are inherited uniformly
from the finite-$k$ systems
(Lem.~\ref{lem:synth-bound})
via Cor.~\ref{cor:proj-strong}.
\end{proof}

\begin{remark}[Bari--Krein, demoted]
\label{rem:bari-role}
Since $\{\Phi_n^{(\infty)}\}$ is orthonormal in
$\mathcal{K}$, it is a Riesz basis with $A=B=1$
by direct computation.
Bari--Krein~\cite{GohbergKrein1969} confirms
this via Gram$=I$, but is not the primary argument.
The non-trivial content is in
Lem.~\ref{lem:proj-conv} (Kato stability),
not in the Gram computation.
\end{remark}

%%════════════════════════════════════════════════════════════════
\section{Canonicalization}
\label{sec:canonical}
%%════════════════════════════════════════════════════════════════

\begin{lemma}[Spectral stability of $\mathcal{K}$]
\label{lem:same-K}
Let $c_k$ and $c_k'$ be two diagonal subsequences
both satisfying Lem.~\ref{lem:XI-diagonal}
and \hyperref[hyp:gap-structural]{Hyp.~Gap-S}.
Then their closed spans $\mathcal{K}$ and $\mathcal{K}'$
coincide.
\end{lemma}
\begin{proof}
Both are closed invariant subspaces of $\Hlim$
(the SOT limit is uniquely determined,
Paper~X Thread~A).
By Lem.~\ref{lem:proj-conv},
both $\{\Pninf\}$ and $\{P_n^{(\infty)'}\}$
are obtained as norm limits of spectral
projections of $\Hlim$.
Since $\Hlim$ has simple spectrum in the transition
regime (\eqref{XII3}+\eqref{XII5} in the limit)
and each spectral projection of $\Hlim$
at an isolated eigenvalue is uniquely determined
(Riesz formula, isolated by Gap-S),
the two families yield the same projections,
hence $\mathcal{K}=\mathcal{K}'$.
\end{proof}

\begin{theorem}[Canonicalization]
\label{thm:canonical-proof}
Under \hyperref[hyp:gap-structural]{Hyp.~Gap-S},
$\Hlim$ is a canonical self-adjoint operator
on $\mathcal{K}$, independent of diagonal subsequence
up to unitary equivalence.
\end{theorem}
\begin{proof}
(1) Riesz basis: Prop.~\ref{prop:riesz-basis}.
(2) Unique $\mathcal{K}$: Lem.~\ref{lem:same-K}.
(3) Spectrum $\{\lambda_n^{(\infty)}\}$
    independent of subsequence
    (SOT uniqueness + \eqref{XII2}).
(4) $\Hlim|_{\mathcal{K}}$ has simple spectrum
    with orthonormal eigenbasis;
    two such operators are unitarily equivalent
    (spectral theorem).
\end{proof}

%%════════════════════════════════════════════════════════════════
\section{System Closure}
\label{sec:closure}
%%════════════════════════════════════════════════════════════════

Under \hyperref[hyp:gap-structural]{Hyp.~Gap-S},
the three open axes of Paper~XI close:
\begin{enumerate}
\item \textbf{Subsequence dependence:}
      Thm.~\ref{thm:canonical-proof}.
\item \textbf{NAcc as axiom:} Cor.~\ref{cor:nacc}.
\item \textbf{Canonical labeling:}
      Riesz-basis indexing of $\{\Phi_n^{(\infty)}\}$
      (Prop.~\ref{prop:riesz-basis}).
\end{enumerate}

\begin{remark}[Remaining open]
Completeness ($\mathcal{K}=L^2(\mathbb{R})$):
Paper~XIII.
Density of $\{\lambda_n^{(\infty)}\}$:
Paper~XI Conj.~6.1.
Sharp gap scale $g(c)$:
Prob.~\ref{prob:sharp-gc}.
\end{remark}

%%════════════════════════════════════════════════════════════════
\section{Open Problems}
\label{sec:open}
%%════════════════════════════════════════════════════════════════

\begin{problem}[Sharp gap scale]
\label{prob:sharp-gc}
Determine the sharp form of $g(c)$.
By Thm.~\ref{thm:wkb-indeterminacy}, this requires
global control of $F_c'$ via:
\begin{enumerate}
\item \textit{Plancherel--Rotach / kernel asymptotics:}
      density of concentration eigenvalues near $\lambda=1/2$.
\item \textit{Resolvent asymptotics:}
      $\operatorname{Im}(\mathcal{H}_c-z)^{-1}$ near $z=1/2+i0$.
\end{enumerate}
Expected: $g(c)=c^{-\alpha}$, $\alpha\in(1/4,1)$;
precise value unknown.
\end{problem}

\begin{problem}[Completeness]
\label{prob:completeness}
Is $\mathcal{K}=L^2(\mathbb{R})$? (Paper~XIII.)
\end{problem}

\begin{problem}[Density]
\label{prob:density}
Is $\{\lambda_n^{(\infty)}\}$ dense in $[0,1]$?
\end{problem}

\begin{problem}[Explicit Riesz constants]
\label{prob:riesz-constants}
Determine $A$, $B$ in Prop.~\ref{prop:riesz-basis}
explicitly in terms of $g(c)$
via the Kato stability bound \eqref{eq:proj-conv}.
\end{problem}

%%%─── Bibliography
\begin{thebibliography}{99}

\bibitem{PaperXI}
[Author],
\textit{Spectral structure and density
for the limiting PSWF operator},
Paper~XI (final), 2026.

\bibitem{PaperX}
[Author],
\textit{Mosco form convergence, resolvent
stability, and the bridge theorem
for PSWF concentration operators},
Paper~X, preprint, 2026.

\bibitem{PaperVIII}
[Author],
\textit{Non-cancellation control, the sharp
$c^{-1/2}$ lower bound, and spectral
connection to the zeros of $\zeta(s)$},
Paper~VIII, preprint, 2026.

\bibitem{Slepian1961}
D.~Slepian and H.~O.~Pollak,
\textit{Prolate spheroidal wave functions,
Fourier analysis and uncertainty~I},
Bell Syst.\ Tech.\ J.\ \textbf{40} (1961), 43--63.

\bibitem{Osipov2013}
A.~Osipov, V.~Rokhlin, and H.~Xiao,
\textit{Prolate Spheroidal Wave Functions
of Order Zero},
Springer, New York, 2013.

\bibitem{Kato1966}
T.~Kato,
\textit{Perturbation Theory for Linear Operators},
Springer, Berlin, 1966.
(Rank-1 projection stability:
\S\,II.1.4, esp.\ Eq.~(II.1.17).)

\bibitem{ReedSimon1975}
M.~Reed and B.~Simon,
\textit{Methods of Modern Mathematical Physics,
Vols.~I--II},
Academic Press, New York, 1972--1975.
(Strong vs.\ norm resolvent convergence:
Vol.~I, Ex.~VIII.7.14.)

\bibitem{GohbergKrein1969}
I.~Gohberg and M.~Krein,
\textit{Introduction to the Theory of
Linear Nonselfadjoint Operators},
AMS, Providence, 1969.

\bibitem{Christensen2003}
O.~Christensen,
\textit{An Introduction to Frames
and Riesz Bases},
Birkh\"auser, Boston, 2003.

\bibitem{Widom1964}
H.~Widom,
\textit{Asymptotic behavior of the
eigenvalues of certain integral
equations~II},
Arch.\ Ration.\ Mech.\ Anal.\ \textbf{17}
(1964), 215--229.

\bibitem{DLMF2023}
F.~W.~J.~Olver et al.\ (eds.),
\textit{NIST Digital Library of Mathematical
Functions}, Release 1.1.10, 2023.
\texttt{https://dlmf.nist.gov/}

\bibitem{CCM2025}
A.~Connes, C.~Consani, and H.~Moscovici,
\textit{Zeta Spectral Triples},
arXiv:2511.22755 [math.NT], 2025.

\end{thebibliography}

\end{document}
